\documentclass[aspectratio=169]{beamer}
\usetheme{Madrid}
\usecolortheme{default}
\usepackage[utf8]{inputenc}
\usepackage[spanish]{babel}
\usepackage{listings}
\usepackage{xcolor}
\usepackage{amsmath}
\usepackage{bookmark}
\usepackage{array}
\usepackage{booktabs}
\usepackage{textcomp}

% ===== ESTILO DE CODIGO =====
\definecolor{backcolour}{rgb}{0.95,0.95,0.92}
\definecolor{codegreen}{rgb}{0,0.5,0}
\definecolor{codegray}{rgb}{0.4,0.4,0.4}
\definecolor{codeblue}{rgb}{0,0,0.7}

\lstdefinestyle{mystyle}{
  backgroundcolor=\color{backcolour},
  basicstyle=\ttfamily\scriptsize,
  keywordstyle=\color{codeblue}\bfseries,
  commentstyle=\color{codegreen},
  stringstyle=\color{red},
  breaklines=true,
  numbers=left,
  numbersep=5pt,
  numberstyle=\tiny\color{codegray},
  language=Python,
  showstringspaces=false,
  frame=single,
  rulecolor=\color{gray!40},
}
\lstset{style=mystyle}

% ===== METADATOS =====
\title{Sistema de Computaci\'on Ubicua}
\subtitle{Comunicaci\'on Bidireccional entre Nodo Central y Nodos Sensores}
\author{Jose Mena P}
\institute{Tecnolog\'ias Emergentes}
\date{\today}

% ===== DOCUMENTO =====
\begin{document}

%------------------------------------------------
\begin{frame}
\titlepage
\end{frame}

%------------------------------------------------
\begin{frame}{Contenido}
\tableofcontents
\end{frame}

%================================================
\section{Computaci\'on Ubicua}
%================================================

\begin{frame}{?`Qu\'e es la Computaci\'on Ubicua?}

La \textbf{computaci\'on ubicua} (ubiquitous computing) es el paradigma donde la
tecnolog\'ia est\'a integrada en el entorno f\'isico de forma transparente para el usuario.

\vspace{0.5em}
\begin{block}{Concepto clave}
Los dispositivos computan y se comunican \textbf{en todo momento y lugar},
sin que el usuario deba interactuar directamente con ellos.
\end{block}

\vspace{0.5em}
\begin{itemize}
  \item Propuesto por Mark Weiser (Xerox PARC, 1991).
  \item Base del Internet de las Cosas (IoT).
  \item Los sensores recolectan datos; el nodo central decide y actua.
  \item Ejemplos: termostatos inteligentes, c\'amaras de seguridad, wearables.
\end{itemize}
\end{frame}

%------------------------------------------------
\begin{frame}{Principios del Sistema Dise\~nado}

Nuestro sistema aplica tres principios fundamentales de la computaci\'on ubicua:

\vspace{0.5em}
\begin{columns}
\column{0.32\textwidth}
\begin{exampleblock}{Percepci\'on}
Los nodos sensores miden temperatura, humedad y movimiento del entorno f\'isico.
\end{exampleblock}

\column{0.32\textwidth}
\begin{exampleblock}{Comunicaci\'on}
Los nodos intercambian informaci\'on con el nodo central mediante el protocolo HTTP.
\end{exampleblock}

\column{0.32\textwidth}
\begin{exampleblock}{Acci\'on}
El nodo central env\'ia comandos para controlar actuadores (LEDs y motores).
\end{exampleblock}
\end{columns}

\vspace{1em}
\begin{alertblock}{Caracter\'istica diferenciadora}
La comunicaci\'on es \textbf{bidireccional}: los nodos no solo env\'ian datos,
tambi\'en reciben y ejecutan \'ordenes del nodo central.
\end{alertblock}
\end{frame}

%================================================
\section{Arquitectura del Sistema}
%================================================

\begin{frame}{Diagrama General del Sistema}

\begin{center}
\ttfamily\small
\begin{tabular}{c}
\hline
\textbf{Usuario (Telegram / Streamlit)} \\
\hline
$\updownarrow$ comandos / alertas \\
\hline
\textbf{NODO CENTRAL} \quad nodo\_central.py \\
Monitoreo + Bot Telegram + Men\'u consola \\
\hline
$\updownarrow$ HTTP GET /datos \quad $\updownarrow$ HTTP POST /comando \\
\hline
\begin{tabular}{c|c}
\textbf{NODO SENSOR 1} & \textbf{NODO SENSOR 2} \\
localhost:5001 & localhost:5002 \\
Temperatura + Humedad & Movimiento + Presencia \\
LED + Motor & LED \\
\end{tabular} \\
\hline
\end{tabular}
\end{center}

\vspace{0.5em}
\begin{itemize}
  \item Toda la comunicaci\'on usa \textbf{HTTP + JSON}.
  \item No requiere librer\'ias externas: solo la biblioteca est\'andar de Python.
  \item Cada nodo corre en su propio proceso (terminal independiente).
\end{itemize}
\end{frame}

%------------------------------------------------
\begin{frame}{Archivos del Proyecto}

\begin{center}
\begin{tabular}{lll}
\toprule
\textbf{Archivo} & \textbf{Puerto} & \textbf{Funci\'on} \\
\midrule
\texttt{config.py}         & ---    & Token Telegram, URLs, umbrales \\
\texttt{nodo\_sensor\_1.py} & 5001   & Servidor HTTP: temperatura, LED, motor \\
\texttt{nodo\_sensor\_2.py} & 5002   & Servidor HTTP: movimiento, LED \\
\texttt{nodo\_central.py}  & ---    & Monitoreo, alertas, bot Telegram \\
\texttt{interfaz.py}       & 8501   & Panel Streamlit (visualizaci\'on) \\
\bottomrule
\end{tabular}
\end{center}

\vspace{0.5em}
\begin{block}{Dependencias}
\texttt{pip install streamlit} \quad (el resto usa solo la biblioteca est\'andar de Python 3)
\end{block}
\end{frame}

%================================================
\section{Nodo Sensor 1}
%================================================

\begin{frame}[fragile]{Nodo Sensor 1 --- Sensores Simulados}

El Nodo 1 simula dos sensores f\'isicos mediante funciones que generan valores aleatorios
con comportamiento realista.

\vspace{0.3em}
\begin{lstlisting}
def leer_temperatura():
    global _temperatura_base
    # Fluctuacion lenta del ambiente
    _temperatura_base += random.uniform(-0.3, 0.3)
    _temperatura_base = max(10.0, min(45.0, _temperatura_base))
    # Pico esporadico (5% probabilidad) para disparar alertas
    if random.random() < 0.05:
        return round(_temperatura_base + random.uniform(10, 15), 1)
    return round(_temperatura_base + random.uniform(-0.5, 0.5), 1)

def leer_humedad():
    return round(random.uniform(40.0, 80.0), 1)
\end{lstlisting}
\end{frame}

%------------------------------------------------
\begin{frame}[fragile]{Nodo Sensor 1 --- Servidor HTTP}

El nodo implementa un servidor HTTP con dos endpoints usando \texttt{http.server}
de la biblioteca est\'andar de Python.

\vspace{0.3em}
\begin{lstlisting}
class ManejadorNodo1(BaseHTTPRequestHandler):

    def do_GET(self):           # GET /datos
        if self.path == "/datos":
            respuesta = {
                "temperatura": leer_temperatura(),
                "humedad":     leer_humedad(),
                "led":         estado["led"],
                "motor":       estado["motor"],
            }
            self._responder_json(200, respuesta)

    def do_POST(self):          # POST /comando
        cuerpo  = json.loads(self.rfile.read(...))
        comando = cuerpo.get("comando", "").upper()
        if   comando == "LED_ON":    estado["led"]   = True
        elif comando == "LED_OFF":   estado["led"]   = False
        elif comando == "MOTOR_ON":  estado["motor"] = True
        elif comando == "MOTOR_OFF": estado["motor"] = False
        self._responder_json(200, {"ok": True, "comando": comando})
\end{lstlisting}
\end{frame}

%================================================
\section{Nodo Sensor 2}
%================================================

\begin{frame}[fragile]{Nodo Sensor 2 --- Sensor de Movimiento}

El Nodo 2 simula un sensor PIR (Passive Infrared) que detecta movimiento.
La probabilidad var\'ia seg\'un la hora del d\'ia para simular actividad real.

\vspace{0.3em}
\begin{lstlisting}
def leer_movimiento():
    hora_actual = datetime.now().hour
    # Mayor probabilidad en horario diurno
    if 8 <= hora_actual <= 18:
        probabilidad = 0.25     # 25% durante el dia
    else:
        probabilidad = 0.10     # 10% en horario nocturno

    detectado = random.random() < probabilidad
    if detectado:
        _detecciones_movimiento += 1
    return detectado
\end{lstlisting}

\vspace{0.3em}
\begin{block}{Actuador}
El Nodo 2 tiene un \textbf{LED de alerta} que el nodo central enciende
autom\'aticamente cuando detecta movimiento fuera del horario permitido.
\end{block}
\end{frame}

%================================================
\section{Nodo Central}
%================================================

\begin{frame}[fragile]{Nodo Central --- Arquitectura de Hilos}

El nodo central ejecuta \textbf{tres tareas en paralelo} usando \texttt{threading}:

\vspace{0.5em}
\begin{columns}
\column{0.32\textwidth}
\begin{block}{Hilo 1\\\small Monitoreo}
Cada 5 segundos consulta ambos nodos, verifica umbrales y registra fallos.
\end{block}

\column{0.32\textwidth}
\begin{block}{Hilo 2\\\small Bot Telegram}
Hace long polling a la API de Telegram y responde comandos del usuario.
\end{block}

\column{0.32\textwidth}
\begin{block}{Hilo Principal\\\small Men\'u consola}
Men\'u interactivo en terminal para enviar comandos manualmente.
\end{block}
\end{columns}

\vspace{1em}
\begin{lstlisting}
hilo_monitoreo = threading.Thread(target=ciclo_monitoreo, daemon=True)
hilo_telegram  = threading.Thread(target=ciclo_bot_telegram, daemon=True)
hilo_monitoreo.start()
hilo_telegram.start()
menu_consola()   # hilo principal (bloquea hasta que el usuario salga)
\end{lstlisting}
\end{frame}

%------------------------------------------------
\begin{frame}[fragile]{Nodo Central --- Ciclo de Monitoreo}

\begin{lstlisting}
def ciclo_monitoreo():
    while True:
        # --- Nodo 1 ---
        datos1 = consultar_nodo(config.NODO_1_URL, "Nodo1")
        if datos1:
            nodo_recuperado("nodo_1", "Nodo1")
            verificar_temperatura(datos1)     # alerta si fuera de umbral
        else:
            verificar_nodo_caido("nodo_1", "Nodo1")  # alerta si sin respuesta

        # --- Nodo 2 ---
        datos2 = consultar_nodo(config.NODO_2_URL, "Nodo2")
        if datos2:
            nodo_recuperado("nodo_2", "Nodo2")
            verificar_movimiento(datos2)      # alerta si movimiento nocturno
        else:
            verificar_nodo_caido("nodo_2", "Nodo2")

        time.sleep(config.INTERVALO_POLLING_SEG)   # esperar 5 segundos
\end{lstlisting}
\end{frame}

%================================================
\section{Bot de Telegram}
%================================================

\begin{frame}[fragile]{Bot de Telegram --- Configuraci\'on}

\begin{columns}
\column{0.55\textwidth}
\textbf{Pasos para crear el bot:}
\begin{enumerate}
  \item Buscar \texttt{@BotFather} en Telegram
  \item Ejecutar \texttt{/newbot} y asignar nombre y username
  \item Copiar el \textbf{token} generado (formato: \texttt{123456789:AABBcc...})
  \item Enviar un mensaje al bot y visitar:\\
        \texttt{api.telegram.org/bot\{TOKEN\}/getUpdates}
  \item Anotar el campo \texttt{"id"} dentro de \texttt{"chat"} (\textbf{Chat ID})
  \item Pegar ambos valores en \texttt{config.py}
\end{enumerate}

\column{0.42\textwidth}
\begin{exampleblock}{config.py}
\begin{lstlisting}
TELEGRAM_TOKEN =
  "8689495203:AAFrKj..."

TELEGRAM_CHAT_ID =
  "933445452"
\end{lstlisting}
\end{exampleblock}
\end{columns}
\end{frame}

%------------------------------------------------
\begin{frame}[fragile]{Bot de Telegram --- Long Polling}

El bot utiliza la t\'ecnica de \textbf{long polling} para recibir mensajes
sin necesidad de un servidor p\'ublico (no requiere webhook).

\vspace{0.3em}
\begin{lstlisting}
def ciclo_bot_telegram():
    # Saltar mensajes previos al inicio del programa
    pendientes = obtener_actualizaciones_telegram(offset=0)
    offset = pendientes[-1]["update_id"] + 1 if pendientes else 0

    while True:
        actualizaciones = obtener_actualizaciones_telegram(offset)
        for update in actualizaciones:
            offset   = update["update_id"] + 1
            texto    = update["message"]["text"]
            chat_id  = str(update["message"]["chat"]["id"])
            respuesta = procesar_comando_telegram(texto)
            # Responder al usuario
            _llamar_api_telegram("sendMessage", {
                "chat_id": chat_id,
                "text":    respuesta,
            })
\end{lstlisting}
\end{frame}

%------------------------------------------------
\begin{frame}{Bot de Telegram --- Comandos Disponibles}

\begin{center}
\begin{tabular}{ll}
\toprule
\textbf{Comando} & \textbf{Acci\'on} \\
\midrule
\texttt{/ayuda}    & Lista todos los comandos disponibles \\
\texttt{/estado}   & Muestra temperatura, movimiento y actuadores \\
\texttt{/led1\_on}  & Enciende el LED del Nodo 1 \\
\texttt{/led1\_off} & Apaga el LED del Nodo 1 \\
\texttt{/motor\_on} & Enciende el motor del Nodo 1 \\
\texttt{/motor\_off}& Apaga el motor del Nodo 1 \\
\texttt{/led2\_on}  & Enciende el LED del Nodo 2 \\
\texttt{/led2\_off} & Apaga el LED del Nodo 2 \\
\bottomrule
\end{tabular}
\end{center}

\vspace{0.5em}
\begin{alertblock}{Alertas autom\'aticas (sin comando del usuario)}
El bot tambi\'en env\'ia mensajes proactivos cuando detecta condiciones
anormales (temperatura, movimiento nocturno, nodo ca\'ido).
\end{alertblock}
\end{frame}

%================================================
\section{Interfaz Streamlit}
%================================================

\begin{frame}{Interfaz Streamlit --- Panel de Control}

La interfaz \texttt{interfaz.py} provee un panel visual que:

\begin{itemize}
  \item Consulta los nodos cada \textbf{3 segundos} (\texttt{st.rerun()})
  \item Muestra el estado de cada actuador con indicadores de color:
    \begin{itemize}
      \item \textcolor{green}{\textbullet} \textbf{Verde} = LED Nodo 1 encendido
      \item \textcolor{blue}{\textbullet} \textbf{Azul} = Motor encendido
      \item \textcolor{orange}{\textbullet} \textbf{Amarillo} = LED Nodo 2 encendido
      \item \textbullet{} \textbf{Negro} = apagado
    \end{itemize}
  \item Incluye botones \textit{Encender / Apagar} para cada actuador
  \item Muestra alertas cuando hay movimiento detectado
\end{itemize}

\vspace{0.5em}
\begin{block}{Ejecutar}
\texttt{streamlit run interfaz.py}\\
Se abre en \texttt{http://localhost:8501}
\end{block}
\end{frame}

%------------------------------------------------
\begin{frame}[fragile]{Interfaz Streamlit --- L\'ogica de Auto-refresco}

\begin{lstlisting}
# Consultar datos de ambos nodos en cada ciclo
datos1 = consultar_nodo(config.NODO_1_URL)
datos2 = consultar_nodo(config.NODO_2_URL)

# Mostrar indicador visual del LED
led1_on = datos1["led"]
color_led1 = "verde" if led1_on else "negro"
st.markdown(f"### {'Encendido' if led1_on else 'Apagado'}")

# Botones de control
if st.button("Encender LED"):
    enviar_comando(config.NODO_1_URL, "LED_ON")
    st.rerun()        # refrescar inmediatamente

# Auto-refresh automatico cada 3 segundos
time.sleep(3)
st.rerun()
\end{lstlisting}
\end{frame}

%================================================
\section{Alertas y Umbrales}
%================================================

\begin{frame}[fragile]{L\'ogica de Alertas --- Temperatura}

\begin{lstlisting}
# Umbrales en config.py
TEMP_MAX_C = 30.0    # Alerta si temperatura SUPERA este valor
TEMP_MIN_C = 15.0    # Alerta si temperatura es INFERIOR a este valor

def verificar_temperatura(datos_nodo1):
    temp = datos_nodo1.get("temperatura")
    if temp > config.TEMP_MAX_C and not _alerta_temp_enviada["alta"]:
        msg = (f"ALERTA TEMPERATURA ALTA\n"
               f"Temperatura: {temp} grados C\n"
               f"Umbral maximo: {config.TEMP_MAX_C} grados C")
        enviar_mensaje_telegram(msg)
        _alerta_temp_enviada["alta"] = True
    elif temp <= config.TEMP_MAX_C:
        _alerta_temp_enviada["alta"] = False   # Resetear alerta
\end{lstlisting}

\begin{block}{Mecanismo anti-spam}
La alerta se env\'ia \textbf{una sola vez} cuando se supera el umbral.
Solo se vuelve a enviar si la temperatura regresa al rango normal y vuelve a subir.
\end{block}
\end{frame}

%------------------------------------------------
\begin{frame}[fragile]{L\'ogica de Alertas --- Movimiento y Nodo Ca\'ido}

\begin{columns}
\column{0.5\textwidth}
\textbf{Movimiento nocturno}
\begin{lstlisting}
# Horario de alerta (config.py)
HORA_ALERTA_INICIO = 22
HORA_ALERTA_FIN    = 6

def verificar_movimiento(datos):
    if not datos.get("movimiento"):
        return
    hora = datetime.now().hour
    fuera = (hora >= 22) or (hora < 6)
    if fuera:
        enviar_mensaje_telegram(
          "ALERTA MOVIMIENTO "
          "NOCTURNO")
\end{lstlisting}

\column{0.5\textwidth}
\textbf{Nodo sin respuesta}
\begin{lstlisting}
MAX_SIN_RESPUESTA = 3

def verificar_nodo_caido(
        nombre_clave,
        nombre_display):
    _fallos[nombre_clave] += 1
    fallos = _fallos[nombre_clave]
    if fallos == MAX_SIN_RESPUESTA:
        enviar_mensaje_telegram(
          f"ALERTA NODO SIN "
          f"RESPUESTA: "
          f"{nombre_display}\n"
          f"Intentos: {fallos}")
\end{lstlisting}
\end{columns}
\end{frame}

%------------------------------------------------
\begin{frame}{Resumen de Umbrales Configurados}

\begin{center}
\begin{tabular}{llll}
\toprule
\textbf{Condici\'on} & \textbf{Umbral} & \textbf{Par\'ametro} & \textbf{Alerta} \\
\midrule
Temperatura alta     & $> 30.0$\,\textdegree C & \texttt{TEMP\_MAX\_C}           & S\'i, \'unica vez \\
Temperatura baja     & $< 15.0$\,\textdegree C & \texttt{TEMP\_MIN\_C}           & S\'i, \'unica vez \\
Movimiento nocturno  & 22:00 -- 06:00           & \texttt{HORA\_ALERTA\_*}        & S\'i, por evento \\
Nodo sin respuesta   & 3 fallos seguidos        & \texttt{MAX\_SIN\_RESPUESTA}    & S\'i, al llegar a 3 \\
\bottomrule
\end{tabular}
\end{center}

\vspace{0.5em}
Todos los umbrales se modifican en \texttt{config.py} sin tocar el c\'odigo principal.
\end{frame}

%================================================
\section{Configuraci\'on Segura}
%================================================

\begin{frame}[fragile]{Gesti\'on de Secretos --- secretos.py}

El token de Telegram \textbf{nunca debe subirse a git}. El sistema usa un archivo
separado excluido del repositorio:

\vspace{0.3em}
\begin{columns}
\column{0.48\textwidth}
\begin{exampleblock}{secretos.example.py \quad (en git)}
\begin{lstlisting}
# Plantilla para estudiantes
TELEGRAM_TOKEN   =
  "PEGA_TU_TOKEN"
TELEGRAM_CHAT_ID =
  "PEGA_TU_CHAT_ID"
\end{lstlisting}
\end{exampleblock}

\column{0.48\textwidth}
\begin{alertblock}{secretos.py \quad (NO en git)}
\begin{lstlisting}
# Valores reales del estudiante
TELEGRAM_TOKEN   =
  "8689495203:AAFrKj..."
TELEGRAM_CHAT_ID =
  "933445452"
\end{lstlisting}
\end{alertblock}
\end{columns}

\vspace{0.5em}
\begin{block}{config.py carga los secretos as\'i}
\begin{lstlisting}
import secretos
TELEGRAM_TOKEN   = secretos.TELEGRAM_TOKEN
TELEGRAM_CHAT_ID = secretos.TELEGRAM_CHAT_ID
\end{lstlisting}
\end{block}
\end{frame}

%------------------------------------------------
\begin{frame}[fragile]{Utilidad obtener\_chat\_id.py}

Script de configuraci\'on inicial: espera un mensaje del estudiante en Telegram
y muestra su Chat ID autom\'aticamente.

\vspace{0.3em}
\begin{lstlisting}
# 1. Saltar mensajes viejos para no procesar el historial
updates = get_updates(timeout=1)
offset = (updates[-1]["update_id"] + 1) if updates else 0

print("Esperando mensaje en Telegram...")

# 2. Long polling: esperar hasta que llegue un mensaje nuevo
while True:
    resp = get_updates(offset=offset, timeout=20)
    for update in resp.get("result", []):
        chat_id = update["message"]["chat"]["id"]
        print(f"Tu Chat ID es: {chat_id}")
        exit(0)   # cerrar al obtener el dato
\end{lstlisting}

\begin{block}{Flujo del estudiante}
\texttt{python obtener\_chat\_id.py}
$\Rightarrow$ escribir al bot en Telegram
$\Rightarrow$ copiar el n\'umero que aparece
$\Rightarrow$ pegarlo en \texttt{secretos.py}
\end{block}
\end{frame}

%================================================
\section{Ejecuci\'on y Verificaci\'on}
%================================================

\begin{frame}[fragile]{Paso 1 --- Configurar Credenciales (una sola vez)}

Antes de correr cualquier archivo, configura tu bot de Telegram:

\vspace{0.4em}
\begin{alertblock}{Accion requerida}
Copia \texttt{secretos.example.py} como \texttt{secretos.py} y pega tu token de BotFather.
\end{alertblock}

\vspace{0.4em}
\begin{lstlisting}
# secretos.py  (NO subir a git)
TELEGRAM_TOKEN   = "PEGA_AQUI_TU_TOKEN"
TELEGRAM_CHAT_ID = ""   # se completa en el siguiente paso
\end{lstlisting}

\vspace{0.4em}
Luego corre \textbf{una sola vez} este script para obtener tu Chat ID:

\begin{block}{Terminal --- Obtener Chat ID}
\texttt{cd "C:\textbackslash{}Users\textbackslash{}...\textbackslash{}SistemaUbicuo"}\\
\texttt{python obtener\_chat\_id.py}\\
{\small $\Rightarrow$ Escribele al bot en Telegram $\Rightarrow$ copia el numero que aparece $\Rightarrow$ pegalo en \texttt{secretos.py}}
\end{block}
\end{frame}

%------------------------------------------------
\begin{frame}{Paso 2 --- Levantar los Nodos (en orden)}

Abre \textbf{4 terminales} y ejecuta cada comando. Respeta el orden:

\vspace{0.3em}
\begin{exampleblock}{Terminal 1 --- Nodo Sensor 1 \quad (primero)}
\texttt{cd "C:\textbackslash{}Users\textbackslash{}...\textbackslash{}SistemaUbicuo"}\\
\texttt{python nodo\_sensor\_1.py}
\end{exampleblock}

\begin{exampleblock}{Terminal 2 --- Nodo Sensor 2 \quad (segundo)}
\texttt{cd "C:\textbackslash{}Users\textbackslash{}...\textbackslash{}SistemaUbicuo"}\\
\texttt{python nodo\_sensor\_2.py}
\end{exampleblock}

\begin{block}{Terminal 3 --- Nodo Central \quad (tercero, cuando los nodos ya corren)}
\texttt{cd "C:\textbackslash{}Users\textbackslash{}...\textbackslash{}SistemaUbicuo"}\\
\texttt{python nodo\_central.py}
\end{block}

\begin{block}{Terminal 4 --- Interfaz Streamlit \quad (opcional)}
\texttt{cd "C:\textbackslash{}Users\textbackslash{}...\textbackslash{}SistemaUbicuo"}\\
\texttt{streamlit run interfaz.py}
\end{block}
\end{frame}

%------------------------------------------------
\begin{frame}{Ejecuci\'on del Sistema --- Resumen}

\begin{center}
\begin{tabular}{clll}
\toprule
\textbf{Orden} & \textbf{Comando} & \textbf{Puerto} & \textbf{Cuando} \\
\midrule
1 & \texttt{python obtener\_chat\_id.py} & ---  & Solo la primera vez \\
2 & \texttt{python nodo\_sensor\_1.py}   & 5001 & Antes del central \\
3 & \texttt{python nodo\_sensor\_2.py}   & 5002 & Antes del central \\
4 & \texttt{python nodo\_central.py}     & ---  & Con nodos activos \\
5 & \texttt{streamlit run interfaz.py}   & 8501 & Opcional \\
\bottomrule
\end{tabular}
\end{center}

\vspace{0.5em}
\begin{alertblock}{Importante}
Si el nodo central arranca antes que los sensores, genera alertas falsas de
``nodo ca\'ido''. Siempre levantar primero los nodos sensores.
\end{alertblock}
\end{frame}

%------------------------------------------------
\begin{frame}{Verificaci\'on del Sistema}

\begin{center}
\begin{tabular}{lll}
\toprule
\textbf{Componente} & \textbf{C\'omo verificar} & \textbf{Resultado esperado} \\
\midrule
Nodo 1 activo     & \texttt{GET localhost:5001/datos}     & JSON con temperatura y LED \\
Nodo 2 activo     & \texttt{GET localhost:5002/datos}     & JSON con movimiento \\
Comando HTTP      & \texttt{POST /comando \{LED\_ON\}}   & \texttt{\{"ok": true\}} \\
Bot Telegram      & Escribir \texttt{/estado} al bot      & Respuesta con datos \\
Alerta Telegram   & Esperar pico de temperatura           & Mensaje autom\'atico \\
Streamlit         & \texttt{localhost:8501}               & Panel con botones \\
\bottomrule
\end{tabular}
\end{center}
\end{frame}

%================================================
\section{Conclusiones}
%================================================

\begin{frame}{Conclusiones}

\begin{itemize}
  \item Se implement\'o un sistema de \textbf{computaci\'on ubicua} completamente funcional
        usando solo Python y la biblioteca est\'andar.

  \vspace{0.4em}
  \item La comunicaci\'on es \textbf{bidireccional}: el nodo central consulta datos
        \textit{y} env\'ia comandos a los nodos sensores.

  \vspace{0.4em}
  \item La integraci\'on con \textbf{Telegram} permite monitoreo y control remoto
        en tiempo real desde cualquier dispositivo.

  \vspace{0.4em}
  \item El sistema detecta y notifica autom\'aticamente \textbf{tres tipos de alertas}:
        temperatura fuera de rango, movimiento nocturno y nodo ca\'ido.

  \vspace{0.4em}
  \item La interfaz \textbf{Streamlit} proporciona una visualizaci\'on clara del
        estado de los actuadores con control por clic.

  \vspace{0.4em}
  \item La arquitectura modular permite escalar f\'acilmente agregando m\'as nodos
        o sensores con m\'inimos cambios en \texttt{config.py}.
\end{itemize}
\end{frame}

%------------------------------------------------
\begin{frame}{}
\begin{center}
{\Huge ?`Preguntas?}\\[1em]
{\large Sistema de Computaci\'on Ubicua}\\[0.5em]
{\normalsize Jose Mena P --- Tecnolog\'ias Emergentes}
\end{center}
\end{frame}

\end{document}
